\chapter{Medicines and addiction}

\section*{Definition and Overview}
Medicines and addiction, also referred to as prescription drug addiction or iatrogenic substance use disorder, is a chronic, relapsing brain disease characterised by the compulsive seeking and use of prescribed pharmaceutical agents despite severe harmful consequences. It involves both physiological dependence (defined by tolerance and withdrawal) and psychological addiction (marked by cravings, loss of control, and compulsive use). The primary classes of medications associated with addiction include prescription opioids (prescribed for pain), benzodiazepines and Z-drugs (prescribed for anxiety and insomnia), and central nervous system stimulants (prescribed for attention deficit hyperactivity disorder). While these medications are highly effective therapeutic tools when used under strict medical supervision, their pharmacodynamics pose a significant risk of misuse and dependency, necessitating careful patient selection, monitoring, and structured management.

\section*{Epidemiology}
Prescription drug addiction is a major global public health crisis. In Australia, prescription drug misuse has surpassed the use of many illicit drugs, with pharmaceutical opioids and benzodiazepines contributing to more accidental overdose deaths than heroin. According to national drug strategy surveys, approximately 4--5\% of Australians report misusing a pharmaceutical drug annually. The demographic profile is broad, affecting individuals across all age groups and socioeconomic backgrounds, and frequently including patients who initially began taking the medication for a legitimate medical condition (such as chronic post-surgical pain or severe generalized anxiety). Middle-aged and older adults are increasingly affected due to higher rates of chronic pain, age-related sleep disturbances, and cumulative polypharmacy.

\section*{Aetiology and Pathophysiology}
The development of prescription drug addiction involves complex neurobiological adaptations, patient-specific vulnerabilities, and healthcare system factors:
\begin{itemize}
    \item \textbf{Neurobiology of reward:} Addictive medications increase dopamine levels within the mesolimbic pathway (the brain's reward system). Repeated stimulation leads to neuroplastic changes, down-regulation of endogenous neurotransmitter receptors, and a resetting of reward thresholds.
    \item \textbf{Opioid receptor binding:} Chronic administration of agonists to mu-opioid receptors (e.g. oxycodone, fentanyl, codeine) leads to receptor desensitisation, physical tolerance, withdrawal hyperexcitability, and opioid-induced hyperalgesia (increased sensitivity to pain).
    \item \textbf{GABA-A receptor modulation:} Benzodiazepines (e.g. diazepam, alprazolam) enhance the inhibitory effect of gamma-aminobutyric acid (GABA). Long-term use leads to compensatory down-regulation of GABA-A receptors, resulting in severe physical dependence and a dangerous, seizure-prone withdrawal state.
    \item \textbf{Dopamine and noradrenaline reuptake inhibition:} Stimulant medications (e.g. methylphenidate, dexamphetamine) block the reuptake of dopamine and noradrenaline. When misused or taken via rapid routes (such as crushing and snorting), they trigger rapid, addictive surges in synaptic catecholamines.
    \item \textbf{Co-occurring mental health disorders:} The presence of untreated depression, anxiety, post-traumatic stress disorder (PTSD), or borderline personality disorder increases the likelihood of self-medication and subsequent dependency.
    \item \textbf{Healthcare access and prescribing practices:} Historical over-prescribing, lack of coordination between multiple treating physicians (doctor shopping), and a lack of access to multidisciplinary non-pharmacological pain or mental health services.
\end{itemize}

\section*{Clinical Features}
The clinical presentation of prescription medicine addiction combines behavioural signs of substance-seeking with physiological symptoms of tolerance and withdrawal:
\begin{itemize}
    \item \textbf{Compulsive use and loss of control:} Taking larger doses of the medication or taking it over a longer period than originally intended, along with persistent, unsuccessful attempts to cut down.
    \item \textbf{Drug-seeking behaviours:} Requesting early refills, claiming lost prescriptions, visiting multiple general practitioners or emergency departments (``doctor shopping''), or purchasing medications online or from illicit sources.
    \item \textbf{Tolerance:} A markedly diminished effect with continued use of the same amount of the drug, or the need for progressively larger doses to achieve the same therapeutic or psychoactive effect.
    \item \textbf{Withdrawal syndrome:} Physical and mental distress upon dose reduction or cessation. Opioid withdrawal presents with lacrimation, rhinorrhoea, sweating, dilated pupils, diarrhoea, and severe muscle aches. Benzodiazepine withdrawal presents with severe insomnia, rebound anxiety, muscle tremors, palpitations, and potentially fatal grand mal seizures.
    \item \textbf{Cravings:} An intense, intrusive urge or desire to take the medication, often triggered by stress or environmental cues.
    \item \textbf{Functional decline:} Neglect of occupational, educational, or domestic responsibilities, alongside social withdrawal and abandonment of hobbies in favour of obtaining and using the drug.
\end{itemize}

\section*{Differential Diagnosis}
Clinicians must carefully differentiate between iatrogenic addiction and other clinical phenomena to ensure appropriate management:
\begin{itemize}
    \item \textbf{Pseudoaddiction:} Drug-seeking behaviours (e.g. demanding higher doses or clock-watching) driven solely by under-treated pain. Unlike true addiction, these behaviours resolve entirely once the pain is adequately controlled with appropriate, non-addictive modalities.
    \item \textbf{Physical dependence vs. Addiction:} Differentiating normal physiological adaptation (tolerance and withdrawal), which occurs predictably in patients taking medications like opioids or beta-blockers as prescribed, from the behavioral dysfunction and compulsive, harmful use that defines addiction.
    \item \textbf{Primary anxiety disorder vs. Benzodiazepine withdrawal:} Distinguishing chronic, baseline anxiety from the acute, severe rebound anxiety experienced during benzodiazepine dose reduction.
    \item \textbf{Illicit substance use disorder:} Determining if the misuse of prescription drugs is an isolated issue or part of a wider polysubstance use disorder involving illicit drugs like heroin or methamphetamine.
\end{itemize}

\section*{When to Seek Medical Care}
Patients or family members concerned about medication dependence, escalation of doses, or difficulty stopping a medication should consult a general practitioner or an addiction specialist.

\textbf{Call 000 (or local emergency services) for:}
\begin{itemize}
    \item Signs of an acute opioid overdose: pin-point pupils, slow or shallow breathing (respiratory depression), blue lips or fingernails, cold/clammy skin, and unresponsiveness.
    \item Seizures, severe confusion, extreme agitation, or visual hallucinations (indicating severe benzodiazepine withdrawal).
    \item Sudden loss of consciousness or cardiac arrest.
    \item Suicidal ideation with intent or self-harm attempts.
\end{itemize}

\section*{Diagnosis and Investigations}
The diagnosis is primarily clinical, supported by prescription registries and laboratory investigations:
\begin{itemize}
    \item \textbf{DSM-5 diagnostic criteria:} Structured evaluation against the 11 diagnostic criteria for a substance use disorder, categorised as mild, moderate, or severe.
    \item \textbf{Real-time prescription monitoring (RTPM):} Consultation of national or state databases (such as SafeScript in Australia) to review the patient's prescribing history, identifying multiple prescribers, rapid dispensing, or high-risk drug combinations.
    \item \textbf{Urine toxicology screening:} Assays to confirm the presence of the prescribed drug, detect non-prescribed substances, and verify that the patient is taking, rather than diverting (selling), the medication.
    \item \textbf{Liver and renal function tests:} Essential to monitor organ health, particularly if the patient is taking high doses of combination products (such as codeine-paracetamol, which carries a high risk of hepatotoxicity).
    \item \textbf{Mental health screening:} Using validated tools (e.g. K10, GAD-7, PHQ-9) to identify and diagnose co-occurring mental health conditions.
\end{itemize}

\section*{Management}
Management must be compassionate, multidisciplinary, and focused on gradual reduction and rehabilitation:
\begin{itemize}
    \item \textbf{Supervised medical taper:} Gradually reducing the medication dose to minimise withdrawal symptoms. For benzodiazepines, this often involves converting a short-acting agent to an equivalent dose of a long-acting agent (like diazepam) before commencing a slow taper.
    \item \textbf{Pharmacotherapy for opioid use disorder (Opioid Agonist Treatment):} Substitution therapy with long-acting opioid agonists or partial agonists, such as buprenorphine (often formulated with naloxone to prevent misuse) or methadone.
    \item \textbf{Psychological therapy:} Cognitive Behavioural Therapy (CBT) to address maladaptive coping strategies, Motivational Interviewing (MI) to build readiness for change, and relapse prevention training.
    \item \textbf{Multidisciplinary pain management:} For patients with chronic pain, incorporating physiotherapy, exercise physiology, mindfulness, and non-addictive adjuvant medications (e.g. gabapentinoids, SNRIs, or simple analgesics).
    \item \textbf{Take-home naloxone programmes:} Providing intranasal or intramuscular naloxone kits to patients and carers to enable rapid reversal of accidental opioid overdoses.
    \item \textbf{Peer support groups:} Facilitating engagement with recovery groups such as Narcotics Anonymous (NA) or SMART Recovery.
\end{itemize}

\section*{Complications}
The complications of prescription drug addiction affect physical health, mental wellbeing, and social stability:
\begin{itemize}
    \item \textbf{Accidental fatal overdose:} Severe respiratory depression leading to brain hypoxia, coma, and death, particularly when combining opioids with benzodiazepines or alcohol.
    \item \textbf{Opioid-induced hyperalgesia:} A paradoxical state where chronic opioid use lowers the pain threshold, worsening the patient's baseline pain.
    \item \textbf{Cognitive and psychomotor impairment:} Chronic drowsiness, memory deficits, poor coordination, and an increased risk of motor vehicle accidents and falls.
    \item \textbf{Severe mental health decline:} Worsening depression, emotional lability, social isolation, and an elevated risk of suicide.
    \item \textbf{Social and economic collapse:} Job loss, financial hardship, legal problems, and the breakdown of familial relationships.
    \item \textbf{Transition to illicit drug use:} In some cases, restriction of prescription opioids can lead patients to purchase illicit street drugs, such as heroin or illicitly manufactured fentanyl.
\end{itemize}

\section*{Prognosis and Outcomes}
Prescription drug addiction is a chronic disease, and the prognosis is highly variable. While relapse is common, long-term recovery is achievable with a coordinated, multidisciplinary treatment plan. Patients engaged in opioid agonist treatment (such as buprenorphine) show significant reductions in mortality, illicit drug use, and social instability. The prognosis is greatly enhanced when pharmacological interventions are combined with long-term psychological support, active management of chronic pain or psychiatric comorbidities, and a strong personal support network.

\section*{Prevention and Self-Care}
Preventative strategies and self-care measures are essential to minimise the risk of developing dependency:
\begin{itemize}
    \item \textbf{Informed discussion with prescribers:} Discuss the addictive potential, side effects, and long-term plans for any new opioid, benzodiazepine, or stimulant before starting treatment.
    \item \textbf{Adhere strictly to directions:} Take medications exactly as prescribed; never increase the dose or frequency without consulting your doctor.
    \item \textbf{Coordinate through a single clinic:} Obtain all prescriptions from a single general practitioner and a single pharmacy to prevent therapeutic duplication or dangerous drug interactions.
    \item \textbf{Secure storage and safe disposal:} Store potentially addictive medications in a locked cabinet out of sight of children and visitors, and return unused medications to a pharmacy for safe, ecological disposal.
    \item \textbf{Utilise non-pharmacological therapies:} Actively participate in alternative therapies for pain and stress, such as physiotherapy, cognitive behavioural therapy for insomnia (CBT-I), and daily mindfulness.
\end{itemize}

\section*{Sources and Further Reading}
The following reputable organisations provide further information on medicines and addiction:
\begin{itemize}
    \item Alcohol and Drug Foundation (ADF). \textit{Information on prescription drug misuse, treatment, and support services.} https://adf.org.au/
    \item National Institute on Drug Abuse (NIDA). \textit{Scientific research and patient guides on prescription drug abuse.} https://nida.nih.gov/
    \item NHS. \textit{Overview of addiction to medicines and coping strategies.} https://www.nhs.uk/
    \item Mayo Clinic. \textit{Prescription drug abuse: prevention, symptoms, and medical treatment.} https://www.mayoclinic.org/
    \item ScriptWise. \textit{Australian advocacy organization for prescription medication safety and dependency support.} https://www.scriptwise.org.au/
\end{itemize}
